
\begin{tikzpicture}[
  scale=0.62,
  well/.style      = {thick, black, line cap=round},
  ipline/.style    = {gray!70, thin},
  photonarr/.style = {photon,   line width=1.0pt, -{Stealth[length=5pt]}},
  elecarr/.style   = {electron, line width=1.1pt, -{Stealth[length=5pt]}},
  tick/.style      = {photon,   line width=1.0pt},
  panellabel/.style= {font=\small},
  regime/.style    = {font=\scriptsize, align=center, text width=2.8cm}
]

% étiquette -Ip (une seule fois, à gauche)
\node[font=\small] at (-2.65,-1) {$-I_p$};

% =====================================================================
% a) SINGLE-PHOTON  (xshift 0) : puits symétrique, flèche pleine
% Paramètre a = 0
% =====================================================================
\begin{scope}[xshift=0cm]
  \clip (-2.3,-2.05) rectangle (2.5,1.7);
  \draw[well] plot[domain=-2.5:-0.1, samples=50] (\x, {-1/(5*\x*\x) + 1});
  \draw[well] plot[domain=0.1:2.5, samples=50] (\x, {-1/(5*\x*\x) + 1});
  \draw[ipline] (-2.2,-1)--(2.2,-1);
  \draw[photonarr] (0,-1)--(0,1);
\end{scope}
\node[panellabel] at (-2.0,2.05) {a)};
\node[regime]     at (0,-2.75)   {single-photon\\ionization};

% =====================================================================
% b) ABOVE-THRESHOLD MPI  (x-shift 5) : même puits, flèche pointillée + rungs
% Paramètre a = 0
% =====================================================================
\begin{scope}[xshift=5cm]
  \clip (-2.3,-2.05) rectangle (2.5,1.7);
  \draw[well] plot[domain=-2.5:-0.1, samples=50] (\x, {-1/(5*\x*\x) + 1});
  \draw[well] plot[domain=0.1:2.5, samples=50] (\x, {-1/(5*\x*\x) + 1});
  \draw[ipline] (-2.2,-1)--(2.2,-1);
  % flèche pointillée
  \draw[photonarr] (0,-1.0)--(0,-0.5);
  \draw[photonarr] (0,-0.5)--(0,-0.0);
  \draw[photonarr] (0,0)--(0,0.5);
  \draw[photonarr] (0,0.5)--(0,1.0);
  

\end{scope}
\node[panellabel] at (3.0,2.05) {b)};
\node[regime]     at (5,-2.75)  {above-threshold\\multi-photon ionization};

% =====================================================================
% c) NON-ADIABATIC TUNNEL  (xshift 10) : barrière large, montée + tunnel
% Paramètre a = 0.1
% =====================================================================
\begin{scope}[xshift=10cm]
  \clip (-2.3,-2.05) rectangle (2.5,1.7);
  \draw[well] plot[domain=-2.5:-0.1, samples=50] (\x, {-1/(5*\x*\x) + 1 - 0.5*\x});
  \draw[well] plot[domain=0.1:2.5, samples=50] (\x, {-1/(5*\x*\x) + 1 - 0.5*\x});
  \draw[ipline] (-2.2,-1)--(2.2,-1);
  % montée multiphotonique (courte, jaune)
  \draw[photonarr] (0,-1.0)--(0,-0.5);
  \draw[photonarr] (0,-0.5)--(0,-0.0);
  % tunnel à travers la barrière (vert, sortie à droite)
  \draw[elecarr] (0,-0)--(2.60,-0);
\end{scope}
\node[panellabel] at (8.0,2.05)  {c)};
\node[regime]     at (10,-2.75)  {non-adiabatic\\tunnel ionization};

% =====================================================================
% d) ADIABATIC TUNNEL  (xshift 15) : barrière plus inclinée, tunnel à -Ip
% Paramètre a = 0.9
% =====================================================================
\begin{scope}[xshift=15cm]
  \clip (-2.3,-2.05) rectangle (2.5,1.7);
  \draw[well] plot[domain=-2.5:-0.1, samples=50] (\x, {-1/(5*\x*\x) + 1 - 1.9*\x});
  \draw[well] plot[domain=0.1:2.5, samples=50] (\x, {-1/(5*\x*\x) + 1 - 1.9*\x});
  \draw[ipline] (-2.2,-1)--(2.2,-1);
  % tunnel horizontal au niveau -Ip
  \draw[elecarr] (0,-1)--(2.25,-1);
\end{scope}
\node[panellabel] at (13.0,2.05) {d)};
\node[regime]     at (15,-2.75)  {adiabatic\\tunnel ionization};

% =====================================================================
% e) OVER-BARRIER  (xshift 20) : barrière abaissée sous -Ip, sortie libre
% Paramètre a = 1.6
% =====================================================================
\begin{scope}[xshift=20cm]
  \clip (-2.3,-2.05) rectangle (2.5,1.7);
  \draw[well] plot[domain=-2.5:-0.1, samples=50] (\x, {-1/(5*\x*\x) + 1 - 2.8*\x});
  \draw[well] plot[domain=0.1:2.5, samples=50] (\x, {-1/(5*\x*\x) + 1 - 2.8*\x});
  \draw[ipline] (-2.2,-1)--(2.2,-1);
  % l'électron passe au-dessus de la barrière supprimée
  \draw[elecarr] (0.0,-1)--(2.30,-1);
\end{scope}
\node[panellabel] at (18.0,2.05) {e)};
\node[regime]     at (20,-2.75)  {over-barrier\\ionization};

\end{tikzpicture}
